% GENERATED FILE. Edit canonical lessons, then run npm run generate:books.
% canonical-source-hash: fnv1a64:7b6b9690414a7c4e
% canonical-lessons: HI-C328-rangoli, HI-C328-jagmagana, HI-C328-jhilmilana, HI-C328-patakha, HI-C328-atisbazi

\chapter{A Night of Lights}
\label{ch:hi-a-night-of-lights}
\hlchaptermodality{328}

\begin{chapteropening}
\textbf{By the end of this chapter:} \emph{I can describe a festival night of lights.}

Complete the last lesson of chapter 328: I can describe a festival night of lights.
\end{chapteropening}

\section[\texorpdfstring{\hi{रंगोली} (raṅgolī)}{रंगोली (raṅgolī)}]{\hi{रंगोली} (raṅgolī) — a rangoli, a coloured pattern on the floor}

\label{lesson:HI-C328-rangoli}



\begin{quote}
Before the new one: say the Hindi for an age, then the Hindi for to be ruined.
\end{quote}

\subsection*{You'll want to know: \hi{रंगोली}}
\textbf{\hi{रंगोली}} — \emph{raṅgolī} — \textquotedblleft{}a rangoli, a coloured pattern on the floor\textquotedblright{}.

\begin{quote}
\hi{मीरा} \hi{ने} \hi{आँगन} \hi{में} \hi{रंगोली} \hi{बनाई।}
\end{quote}

\begin{quote}
\emph{Mīrā ne āṁgan meṁ raṅgolī banāī.}
\end{quote}

\begin{quote}
Mira made a rangoli in the courtyard.
\end{quote}

Feminine.

\begin{cousinweb}
\textbf{\hi{रंगोली}} (\emph{raṅgolī}) is probably built on \textbf{\hi{रंग}}, colour, from Sanskrit \emph{raṅga} (Monier-Williams); the rest of the word is not settled. It is drawn on the floor at the door for a festival.
\end{cousinweb}

\begin{grammarlens}[title={What you've built}]
One more word for telling what happened.
\end{grammarlens}

\subsection*{Your turn}
\begin{itemize}
\raggedright
  \item \emph{Say it:} \emph{raṅgolī}
  \item \emph{Say it:} \emph{raṅgolī}, once more
  \item \emph{Say it:} \emph{Mīrā ne āṁgan meṁ raṅgolī banāī.} — Mira made a rangoli in the courtyard.
  \item \emph{Recall:} say the Hindi for an age, then the Hindi for to be ruined, then say \emph{raṅgolī} again
\end{itemize}

\begin{tcolorbox}[breakable,colback=teal!4,colframe=teal!35!black,title={Before you move on}]
What does \hi{रंगोली} mean? (\textquotedblleft{}A rangoli, a coloured pattern on the floor\textquotedblright{}.) Say it once more.
\end{tcolorbox}
\section[\texorpdfstring{\hi{जगमगाना} (jagmagānā)}{जगमगाना (jagmagānā)}]{\hi{जगमगाना} (jagmagānā) — to glitter; to be lit up}

\label{lesson:HI-C328-jagmagana}



\begin{quote}
Before the new one: say the Hindi for to be ruined, then the Hindi for a rangoli, a coloured pattern on the floor.
\end{quote}

\subsection*{You'll want to know: \hi{जगमगाना}}
\textbf{\hi{जगमगाना}} — \emph{jagmagānā} — \textquotedblleft{}to glitter; to be lit up\textquotedblright{}.

\begin{quote}
\hi{त्योहार} \hi{की} \hi{रात} \hi{पूरा} \hi{शहर} \hi{जगमगा} \hi{रहा} \hi{था।}
\end{quote}

\begin{quote}
\emph{tyohār kī rāt pūrā śahar jagmagā rahā thā.}
\end{quote}

\begin{quote}
On the night of the festival the whole city was lit up.
\end{quote}

Past (it …): \hi{वह} \hi{घर} \hi{जगमगाया}, \emph{vah ghar jagmagāyā}.

\begin{cousinweb}
\textbf{\hi{जगमगाना}} (\emph{jagmagānā}) is the verb of \emph{jagmag}, glittering, which CDIAL 5076 sets under \emph{*jag}, a form it reconstructs for shining: a place full of lights \textbf{\hi{जगमगाता} \hi{है}}.
\end{cousinweb}

\begin{grammarlens}[title={What you've built}]
One more word for telling what happened.
\end{grammarlens}

\subsection*{Your turn}
\begin{itemize}
\raggedright
  \item \emph{Say it:} \emph{jagmagānā}
  \item \emph{Say it:} \emph{jagmagānā}, once more
  \item \emph{Say it:} \emph{tyohār kī rāt pūrā śahar jagmagā rahā thā.} — On the night of the festival the whole city was lit up.
  \item \emph{Recall:} say the Hindi for to be ruined, then the Hindi for a rangoli, a coloured pattern on the floor, then say \emph{jagmagānā} again
\end{itemize}

\begin{tcolorbox}[breakable,colback=teal!4,colframe=teal!35!black,title={Before you move on}]
What does \hi{जगमगाना} mean? (\textquotedblleft{}To glitter; to be lit up\textquotedblright{}.) Say it once more.
\end{tcolorbox}
\section[\texorpdfstring{\hi{झिलमिलाना} (jhilmilānā)}{झिलमिलाना (jhilmilānā)}]{\hi{झिलमिलाना} (jhilmilānā) — to twinkle; to flicker}

\label{lesson:HI-C328-jhilmilana}



\begin{quote}
Before the new one: say the Hindi for a rangoli, a coloured pattern on the floor, then the Hindi for to glitter.
\end{quote}

\subsection*{You'll want to know: \hi{झिलमिलाना}}
\textbf{\hi{झिलमिलाना}} — \emph{jhilmilānā} — \textquotedblleft{}to twinkle; to flicker\textquotedblright{}.

\begin{quote}
\hi{खिड़कियों} \hi{पर} \hi{दीये} \hi{झिलमिला} \hi{रहे} \hi{थे।}
\end{quote}

\begin{quote}
\emph{khiṛkiyoṁ par dīye jhilmilā rahe the.}
\end{quote}

\begin{quote}
Lamps were flickering on the windowsills.
\end{quote}

Past (it …): \hi{वह} \hi{तारा} \hi{झिलमिलाया}, \emph{vah tārā jhilmilāyā}.

\begin{cousinweb}
\textbf{\hi{झिलमिलाना}} (\emph{jhilmilānā}) is the verb of \emph{jhilmil}, flickering, which CDIAL 5391 sets under \emph{*jhil}: small lights and stars \textbf{\hi{झिलमिलाते} \hi{हैं}}.
\end{cousinweb}

\begin{grammarlens}[title={What you've built}]
One more word for telling what happened.
\end{grammarlens}

\subsection*{Your turn}
\begin{itemize}
\raggedright
  \item \emph{Say it:} \emph{jhilmilānā}
  \item \emph{Say it:} \emph{jhilmilānā}, once more
  \item \emph{Say it:} \emph{khiṛkiyoṁ par dīye jhilmilā rahe the.} — Lamps were flickering on the windowsills.
  \item \emph{Recall:} say the Hindi for a rangoli, a coloured pattern on the floor, then the Hindi for to glitter, then say \emph{jhilmilānā} again
\end{itemize}

\begin{tcolorbox}[breakable,colback=teal!4,colframe=teal!35!black,title={Before you move on}]
What does \hi{झिलमिलाना} mean? (\textquotedblleft{}To twinkle; to flicker\textquotedblright{}.) Say it once more.
\end{tcolorbox}
\section[\texorpdfstring{\hi{पटाख़ा} (paṭākhā)}{पटाख़ा (paṭākhā)}]{\hi{पटाख़ा} (paṭākhā) — a firecracker}

\label{lesson:HI-C328-patakha}



\begin{quote}
Before the new one: say the Hindi for to glitter, then the Hindi for to twinkle.
\end{quote}

\subsection*{You'll want to know: \hi{पटाख़ा}}
\textbf{\hi{पटाख़ा}} — \emph{paṭākhā} — \textquotedblleft{}a firecracker\textquotedblright{}.

\begin{quote}
\hi{बच्चों} \hi{ने} \hi{रात} \hi{भर} \hi{पटाख़े} \hi{चलाए।}
\end{quote}

\begin{quote}
\emph{baccoṁ ne rāt bhar paṭākhe calāe.}
\end{quote}

\begin{quote}
The children set off firecrackers all night.
\end{quote}

Masculine.

\begin{cousinweb}
\textbf{\hi{पटाख़ा}} (\emph{paṭākhā}) is probably imitative, built on \emph{paṭ}, the sound of a sharp crack. You \textbf{\hi{पटाख़े} \hi{चलाते} \hi{हो}}: you set them off.
\end{cousinweb}

\begin{grammarlens}[title={What you've built}]
One more word for telling what happened.
\end{grammarlens}

\subsection*{Your turn}
\begin{itemize}
\raggedright
  \item \emph{Say it:} \emph{paṭākhā}
  \item \emph{Say it:} \emph{paṭākhā}, once more
  \item \emph{Say it:} \emph{baccoṁ ne rāt bhar paṭākhe calāe.} — The children set off firecrackers all night.
  \item \emph{Recall:} say the Hindi for to glitter, then the Hindi for to twinkle, then say \emph{paṭākhā} again
\end{itemize}

\begin{tcolorbox}[breakable,colback=teal!4,colframe=teal!35!black,title={Before you move on}]
What does \hi{पटाख़ा} mean? (\textquotedblleft{}A firecracker\textquotedblright{}.) Say it once more.
\end{tcolorbox}
\section[\texorpdfstring{\hi{आतिशबाज़ी} (ātiśbāzī)}{आतिशबाज़ी (ātiśbāzī)}]{\hi{आतिशबाज़ी} (ātiśbāzī) — fireworks}

\label{lesson:HI-C328-atisbazi}



\begin{quote}
Before the new one: say the Hindi for to twinkle, then the Hindi for a firecracker.
\end{quote}

\subsection*{You'll want to know: \hi{आतिशबाज़ी}}
\textbf{\hi{आतिशबाज़ी}} — \emph{ātiśbāzī} — \textquotedblleft{}fireworks\textquotedblright{}.

\begin{quote}
\hi{आतिशबाज़ी} \hi{से} \hi{आसमान} \hi{रंगों} \hi{से} \hi{भर} \hi{गया।}
\end{quote}

\begin{quote}
\emph{ātiśbāzī se āsmān raṅgoṁ se bhar gayā.}
\end{quote}

\begin{quote}
The fireworks filled the sky with colour.
\end{quote}

Feminine.

\begin{cousinweb}
\textbf{\hi{आतिशबाज़ी}} (\emph{ātiśbāzī}) is usually traced to Persian \emph{ātišbāzī}, built from \emph{ātiš}, fire, and \emph{bāzī}, play: playing with fire, for a festival.
\end{cousinweb}

\begin{grammarlens}[title={What you've built}]
One more word for telling what happened.
\end{grammarlens}

\subsection*{Your turn}
\begin{itemize}
\raggedright
  \item \emph{Say it:} \emph{ātiśbāzī}
  \item \emph{Say it:} \emph{ātiśbāzī}, once more
  \item \emph{Say it:} \emph{ātiśbāzī se āsmān raṅgoṁ se bhar gayā.} — The fireworks filled the sky with colour.
  \item \emph{Recall:} say the Hindi for to twinkle, then the Hindi for a firecracker, then say \emph{ātiśbāzī} again
\end{itemize}

\begin{tcolorbox}[breakable,colback=teal!4,colframe=teal!35!black,title={Before you move on}]
What does \hi{आतिशबाज़ी} mean? (\textquotedblleft{}Fireworks\textquotedblright{}.) Say it once more.
\end{tcolorbox}
